\begin{abstract}
Sự gia tăng mức độ phức tạp của các cuộc tấn công mạng đặt ra thách thức lớn cho hệ thống phát hiện xâm nhập mạng (NIDS): vừa phải đạt độ chính xác cao trên dữ liệu mất cân bằng, vừa phải tối ưu tài nguyên để triển khai trên các thiết bị điện toán biên (edge computing). Để giải quyết bài toán đa mục tiêu này, nghiên cứu đề xuất kiến trúc NIDS hai giai đoạn dành cho bài toán phân loại nhị phân trên tập dữ liệu CICIDS2017. Giai đoạn một xây dựng một mô hình Giáo viên (Teacher) dựa trên kiến trúc Stacking (kết hợp XGBoost, LightGBM và Random Forest). Để khắc phục mất cân bằng, dữ liệu huấn luyện được tái cân bằng bằng chiến lược lấy mẫu lai RUS (giảm lấy mẫu lớp đa số) và ENN (làm sạch nhiễu ở vùng ranh giới). Giai đoạn hai áp dụng kỹ thuật chưng cất tri thức (Knowledge Distillation) để chuyển giao tri thức từ Teacher sang một mô hình Học sinh (Student) cây quyết định nông. Kết quả thực nghiệm cho thấy mô hình Teacher đạt độ chính xác trên 99,8\%. Đáng chú ý, mô hình Student sau chưng cất vẫn duy trì độ chính xác cao (trên 99\%), đồng thời tăng tốc độ suy luận gấp 368 lần và giảm dung lượng lưu trữ tới 1758 lần. Giải pháp cho thấy tính hiệu quả và khả thi cho các hệ thống an ninh mạng thời gian thực trên thiết bị IoT.
\end{abstract}
